\chapter{Construction~: tick-to-trade, mesuré}\label{ch:ll:build-tick-to-trade-measured}

Un paquet quitte la bourse simulée sur ses deux lignes à la fois. Quelques microsecondes plus tard, le \omterm{def:ll:the-feed-handler:handler}{feed handler} le reçoit, et une microseconde
après, la stratégie a décidé, et un ordre part vers la plateforme en passant par la \omterm{def:ll:pre-trade-risk-and-kill-switches:gate}{barrière de risque} (risk gate) et la \omterm{def:ll:order-gateway-and-order-management:gateway}{passerelle d'ordres}.
Chacune de ces étapes a été construite et mesurée isolément dans ce livre~; ce chapitre les assemble sur un seul chemin, sur quatre threads
d'une même machine, et mesure l'ensemble avec un horodatage à chaque frontière. Le résultat est le \omterm{def:ll:where-latency-comes-from:budget}{budget de latence} du chapitre~1, où les
suppositions sont remplacées par des mesures~: quelle étape coûte combien, à la médiane et dans la queue de distribution, et où un vrai serveur devrait
être différent. La réponse, sur cet ordinateur portable, est sévère sur un point et rassurante sur un autre. Le code écrit dans ce livre
prend quelques centaines de nanosecondes~; le noyau, la pile réseau et une machine partagée prennent le reste, et ce reste en est l'essentiel.

\section{Assemblage du chemin}

Le chemin est constitué des composants du livre, dans l'ordre (\cref{fig:ll:build-tick-to-trade-measured:path}). Un thread \emph{bourse} rejoue
la ligne enregistrée~: la journée à petit pas de cotation du chapitre~19, 6\,000 événements dans 1\,200 paquets MoldUDP64, envoyés en UDP sur l'interface
de bouclage (loopback) sur la ligne A et, cinq microsecondes plus tard, sur la ligne B, à vingt fois son espacement enregistré (un paquet toutes les
\qty{100}{\micro\second}). Un thread \emph{feed} reçoit les deux lignes et exécute le \omterm{def:ll:the-feed-handler:handler}{feed handler} du chapitre~18, qui
les arbitre et publie chaque événement sur un anneau (chapitre~12). Un thread \emph{moteur} prélève chaque événement de l'anneau et exécute
dans \texttt{firm.ticktotrade}, son \texttt{Path}~: le moteur de stratégie du chapitre~20, avec son \omterm{def:ll:the-order-book-builder:builder}{constructeur de carnet} (book builder, chapitre~19), puis, pour chaque
ordre qu'il décide, la \omterm{def:ll:pre-trade-risk-and-kill-switches:gate}{barrière de risque} du chapitre~22 et la \omterm{def:ll:order-gateway-and-order-management:gateway}{passerelle d'ordres} du chapitre~21, qui écrit le message binaire de la plateforme
(chapitre~16) et l'envoie sur une session TCP. Un thread \emph{plateforme} reçoit les trames de la session. Chaque thread est épinglé à son
propre CPU (chapitre~4), et rien sur le chemin n'alloue après la phase de chauffe (chapitre~6)~: le harnais le compte et trouve zéro dans
chaque exécution.

\begin{omfigure}
\begin{tikzpicture}[font=\scriptsize,>=Stealth,box/.style={draw,rounded corners=2pt,minimum height=0.8cm,align=center,
  inner sep=3pt,fill=omDef!12},pt/.style={circle,fill=omThm,inner sep=1.4pt}]
\node[box,fill=gray!12] (ex) at (0,0) {bourse\\ thread};
\node[box] (fd) at (3.3,0) {feed handler\\ (thread feed)};
\node[box,fill=omMeth!15] (rg) at (5.6,0) {anneau};
\node[box,minimum width=4.2cm] (en) at (8.8,0) {thread moteur~:\\ carnet, stratégie, barrière\\ de risque, passerelle};
\node[box,fill=gray!12] (vn) at (12.4,0) {plateforme\\ thread};
\draw[->] (ex) -- node[above=1pt,font=\tiny]{UDP A, B} (fd);
\draw[->] (fd) -- (rg);
\draw[->] (rg) -- (en);
\draw[->] (en) -- node[above]{TCP} (vn);
\foreach \x in {0.75,1.95,4.2,6.15,11.35,11.6} {\node[pt] at (\x,-0.62) {};}
\node[below,align=center] at (0.75,-0.7) {bourse};
\node[below,align=center] at (1.95,-0.7) {reçu};
\node[below,align=center] at (4.2,-0.7) {publ.};
\node[below,align=center] at (6.15,-0.7) {prél.};
\node[below,align=center] at (8.8,-0.7) {carnet, décision,\\ risque, encodage};
\node[below,align=center] at (11.9,-0.7) {envoi,\\ plateforme};
\foreach \x in {7.3,8.2,9.2,10.2} {\node[pt] at (\x,-0.62) {};}
\draw[gray] (0.3,-0.62) -- (12.8,-0.62);
\end{tikzpicture}

\omcaption{Le chemin assemblé et ses \omterm{def:ll:build-tick-to-trade-measured:points}{points d'instrumentation}. Quatre threads sur quatre CPU d'un même ordinateur portable~; chaque ordre porte dix
lectures du \omterm{def:ll:cpu-microarchitecture:tsc}{compteur d'horodatage} (TSC), depuis l'envoi par la bourse du paquet qui l'a provoqué jusqu'à la réception par la
plateforme de l'ordre.}\label{fig:ll:build-tick-to-trade-measured:path}
\end{omfigure}

Le chemin est vérifié avant d'être mesuré. Le test C++ rejoue les mêmes événements hors ligne dans le même \texttt{Path} et
obtient 1\,050 ordres (nouveaux ordres, remplacements et annulations) dont l'empreinte est figée~; le jumeau Rust, assemblé à partir du \omterm{def:ll:the-feed-handler:handler}{feed
handler}, moteur, \omterm{def:ll:pre-trade-risk-and-kill-switches:gate}{barrière de risque} et passerelle, produit la même empreinte~; le harnais en direct, avec ses threads, ses sockets et
l'arbitrage entre les deux lignes, la reproduit à chaque exécution. Ce qui est mesuré est donc le chemin qui calcule les bons
ordres.

\omcode{code/firm/ticktotrade/cpp/firm_ticktotrade.hpp}{82}{91}{Le travail du thread moteur pour un événement de marché~: le moteur
décide, puis chaque ordre traverse la barrière et la passerelle.}\label{lst:ll:build-tick-to-trade-measured:path}

\section{Points d'instrumentation et horodatages des étapes}

\begin{definition}[Point d'instrumentation, attribution de la latence]\label{def:ll:build-tick-to-trade-measured:points}
Un \emph{point d'instrumentation}\index{point d'instrumentation} est un endroit d'un chemin où le temps est lu et enregistré avec
l'identifiant du message qui le traverse. L'\emph{attribution de la latence}\index{attribution de la latence} découpe la latence de bout en bout d'un message
en intervalles consécutifs entre points d'instrumentation, de sorte que chaque partie soit imputée à l'étape qui l'a consommée.
\end{definition}

Chaque point lit le \omterm{def:ll:cpu-microarchitecture:tsc}{compteur d'horodatage} (chapitre~5), que tous les cœurs de cette machine partagent, de sorte qu'une lecture sur le feed
thread et une sur le thread moteur sont comparables. Les dix points de la \cref{fig:ll:build-tick-to-trade-measured:path} donnent huit
étapes~: \emph{réception}, de l'envoi \texttt{sendto} par la bourse jusqu'à l'appel \texttt{recv} du thread feed, soit le chemin UDP du noyau sur
l'interface de loopback~; \emph{décodage}, l'arbitrage et la normalisation par le handler jusqu'à ce que l'événement soit sur l'anneau~; \emph{anneau},
l'attente sur l'anneau~; \emph{carnet}, \emph{stratégie}, \emph{risque} et \emph{passerelle} sur le thread moteur~; et \emph{transmission},
de l'écriture du message à la réception par le thread plateforme, soit le chemin TCP du noyau. Les horodatages voyagent avec l'événement (dans le
message de l'anneau) et avec l'ordre (dans un enregistrement préalloué), et les horodatages de la bourse et de la plateforme sont joints a posteriori par
numéro de séquence et identifiant d'ordre client (\cref{lst:ll:build-tick-to-trade-measured:loop}). Deux pièges ont façonné la conception.
La lecture de l'horloge n'est pas gratuite~: chaque point coûte environ \qty{11}{\nano\second} avec \texttt{rdtscp} (chapitre~5), donc les points
sont peu nombreux et aucun n'est à l'intérieur d'une boucle. Et
lorsqu'un événement provoque deux ordres, le second attend que le premier soit envoyé, quelques microsecondes que l'attribution naïve
impute à son contrôle de risque~; les étapes sont donc attribuées au premier ordre de chaque événement, et les chiffres de bout en bout sur
chaque ordre.

\omcode{code/firm/ticktotrade/cpp/firm_ticktotrade_live.hpp}{162}{187}{Le thread moteur~: prélever un événement de l'anneau, l'horodater,
exécuter le chemin, envoyer chaque ordre et l'horodater.}\label{lst:ll:build-tick-to-trade-measured:loop}

\section{Le budget confronté à la mesure}

La \Cref{fig:ll:build-tick-to-trade-measured:stages} donne les étapes sur dix exécutions, 5\,250 premiers ordres (10\,500 au total), sur cet
ordinateur portable sans cœurs isolés, la machine étant par ailleurs inactive. Le code propre au livre est rapide~: la mise à jour du carnet prend environ
\qty{160}{\nano\second} à la médiane, la stratégie environ 100, la \omterm{def:ll:pre-trade-risk-and-kill-switches:gate}{barrière de risque} environ 50 et la passerelle environ 30, toutes bien en dessous d'une
microseconde au 99\up{e}~centile. L'étape de décodage, environ \qty{440}{\nano\second}, inclut l'attente des messages antérieurs
du même paquet~: cinq événements arrivent ensemble et sortent les uns après les autres. La médiane de l'anneau, environ \qty{0.8}{\micro\second}, est la
même file d'attente, côté moteur, derrière les ordres des événements précédents. Les étapes du noyau écrasent le reste~: la réception d'un
paquet UDP sur l'interface de loopback prend environ \qty{1.4}{\micro\second}, l'envoi de l'ordre en TCP environ \qty{2.9}{\micro\second},
et dans la queue, l'anneau attend des millisecondes chaque fois que le thread moteur est préempté sur son CPU (par le noyau, l'hyperviseur
ou par toute autre tâche),
ce qu'aucun code en espace utilisateur ne peut empêcher sur une machine qui ne l'a pas isolé (chapitre~13).

\begin{omfigure}
\begin{tikzpicture}
\begin{axis}[width=9cm,height=7.2cm,xbar,xmode=log,xmin=10,xmax=5e6,log origin=infty,bar width=3.2pt,
  ytick={0,1,2,3,4,5,6,7,8,9},yticklabels={réception,décodage,anneau,carnet,stratégie,risque,passerelle,transmission,tick-to-trade,en traitement},
  y dir=reverse,xlabel={nanosecondes},tick label style={font=\small},label style={font=\small},grid=major,
  grid style={gray!25},legend pos=outer north east,legend style={font=\scriptsize,cells={anchor=west}},table/col sep=comma]
\addplot[fill=omDef!55,draw=omDef] table[y expr=\coordindex,x=p50_ns]{figdata/low-latency/26-build-tick-to-trade-measured/measured_stages.csv};
\addlegendentry{médiane}
\addplot[fill=omMeth!45,draw=omMeth] table[y expr=\coordindex,x=p99_ns]{figdata/low-latency/26-build-tick-to-trade-measured/measured_stages.csv};
\addlegendentry{99\up{e}~centile}
\addplot[fill=omThm!40,draw=omThm] table[y expr=\coordindex,x=p999_ns]{figdata/low-latency/26-build-tick-to-trade-measured/measured_stages.csv};
\addlegendentry{99.9\up{e}~centile}
\end{axis}
\end{tikzpicture}

\omcaption{Le chemin assemblé, étape par étape et de bout en bout~: médianes, 99\up{e} et 99.9\up{e}~centiles sur dix passes
(les étapes sur le premier ordre de chaque événement, les totaux sur les 10\,500 ordres), sur un ordinateur portable (Intel Core Ultra 7 155H, WSL2) sans cœurs isolés,
la machine étant par ailleurs au repos. Données~: \texttt{bench\_t2t.py}.}\label{fig:ll:build-tick-to-trade-measured:stages}
\end{omfigure}

De bout en bout, de l'envoi du paquet par la bourse à la réception de l'ordre qu'il a provoqué par la plateforme, la médiane vaut environ
\qty{9}{\micro\second}~; dans le processus, de la réception du paquet au message de l'ordre, environ \qty{3.7}{\micro\second}
(la \cref{fig:ll:build-tick-to-trade-measured:ccdf}). Les distributions ont la forme que le livre annonce depuis le chapitre~1~: un corps
serré et une queue située trois ordres de grandeur plus loin, fixée par la machine plutôt que par le code. La \cref{tab:ll:build-tick-to-trade-measured:budget}
vérifie le budget de trois microsecondes du chapitre~1 avec \texttt{firm.latbudget}~: la stratégie et la \omterm{def:ll:pre-trade-risk-and-kill-switches:gate}{barrière de risque} atteignent leurs cibles~;
le \omterm{def:ll:the-order-book-builder:builder}{constructeur de carnet} et le décodage les manquent (le carnet d'un tiers, le décodage à cause de la file d'attente du paquet), et la réception
et la transmission par le noyau les manquent d'un facteur 1.6 et 3.6, comme il fallait s'y attendre~: le budget supposait une
carte réseau à \omterm{def:ll:linux-tuning:polling}{contournement du noyau}.

\begin{omfigure}
\begin{tikzpicture}
\begin{axis}[width=10cm,height=5.8cm,xmode=log,ymode=log,xmin=1000,xmax=5e6,ymin=0.001,ymax=1.1,
  xlabel={nanosecondes},ylabel={fraction des ordres plus lents},tick label style={font=\small},label style={font=\small},
  grid=major,grid style={gray!25},legend pos=south west,legend style={font=\scriptsize,cells={anchor=west}},
  table/col sep=comma]
\addplot[omThm,thick,restrict expr to domain={\coordindex}{0:60}] table[x=ns,y=ccdf]{figdata/low-latency/26-build-tick-to-trade-measured/measured_t2t_ccdf.csv};
\addlegendentry{tick to trade (de la bourse à la plateforme)}
\addplot[omDef,thick,restrict expr to domain={\coordindex}{61:121}] table[x=ns,y=ccdf]{figdata/low-latency/26-build-tick-to-trade-measured/measured_t2t_ccdf.csv};
\addlegendentry{dans le processus (réception au message)}
\end{axis}
\end{tikzpicture}

\omcaption{Les distributions de bout en bout du chemin assemblé, en fraction des 10\,500 ordres plus lents que chaque latence
(échelles logarithmiques), sur le même ordinateur portable. Données~: \texttt{bench\_t2t.py}.}\label{fig:ll:build-tick-to-trade-measured:ccdf}
\end{omfigure}

\begin{table}[ht]
\centering\small
\begin{tabular}{lrrl}
\toprule
étape (médiane) & budget (ns) & mesuré (ns) & \\
\midrule
réception & 900 & 1\,400 & au-dessus \\
décodage & 60 & 440 & au-dessus \\
carnet & 120 & 160 & au-dessus \\
stratégie & 200 & 100 & dans le budget \\
risque & 50 & 47 & dans le budget \\
transmission & 800 & 2\,900 & au-dessus \\
de bout en bout & 3\,000 & 7\,400 & au-dessus \\
\bottomrule
\end{tabular}
\caption{Budget de trois microsecondes du chapitre~1 vérifié par rapport aux étapes mesurées avec \texttt{firm.latbudget} (médianes~; le
bout en bout composé à partir des étapes des premiers ordres). Données~: \texttt{bench\_t2t.py}.}\label{tab:ll:build-tick-to-trade-measured:budget}
\end{table}

\section{Ce qu'un serveur réglé changerait}

Trois changements séparent cet ordinateur portable d'un serveur de trading, et chacun vise une partie du tableau. Une carte réseau à \omterm{def:ll:linux-tuning:polling}{contournement du noyau}
supprime les appels système du chemin des données~: l'interface \texttt{ef\_vi} des adaptateurs d'un fabricant, par exemple, donne à l'application un
«~accès direct~» au chemin de données de l'adaptateur, avec des opérations sur ce chemin qui «~ne nécessitent pas d'appels système~»~;
le budget du chapitre~1 supposait \qty{0.9}{\micro\second} pour recevoir et \qty{0.8}{\micro\second} pour transmettre sur un tel chemin, carte et câble compris. Les cœurs isolés
(chapitre~13) tiennent les autres travaux à l'écart des processeurs du chemin, ce qui constitue la queue de l'étape de l'anneau. Enfin, le chemin lui-même
peut être raccourci~: la file d'attente du paquet derrière le décodage et l'anneau disparaît si les événements sont remis par lots ou si
le fil du flux exécute la stratégie dans les cas simples, et les \qty{160}{\nano\second} du carnet sont l'échelle du chapitre~19 sur un
instrument à petit tick, avec un appel de moteur par événement que le lot amortirait.

Remplacer la réception et la transmission mesurées par les valeurs du budget à \omterm{def:ll:linux-tuning:polling}{contournement du noyau}, et conserver les étapes mesurées dans le processus,
donne une médiane projetée d'environ \qty{3.2}{\micro\second}~: $0.9 + 0.8$ pour le réseau, auxquels s'ajoutent environ \qty{1.5}{\micro\second} pour
le décodage, l'anneau, le carnet, la stratégie, le risque et la passerelle, dont l'attente dans le décodage et l'anneau représente environ les deux tiers. Cela dépasse le
budget de trois microsecondes à cause de la file d'attente, qui est la prochaine chose à supprimer, et non à cause d'un composant que le livre a construit. La projection est
un calcul arithmétique reposant sur une hypothèse énoncée, non une mesure~; la mesure sur le serveur cible est la seule preuve, et le harnais
de ce chapitre sert à la réaliser.

\section{Tutoriel~: faire tourner le chemin, lire le budget}\label{tut:ll:build-tick-to-trade-measured:tut}

\begin{tutorial}
\textbf{Objectif.} Vérifier le chemin assemblé, le faire tourner en direct sur l'interface de bouclage et lire ses étapes par rapport au budget.
\textbf{État final~:} \cref{fig:ll:build-tick-to-trade-measured:stages,fig:ll:build-tick-to-trade-measured:ccdf}, le tableau
du budget, et des tests verts en C++ et en Rust.
\begin{tutsteps}
\item \textbf{La ligne.} \texttt{make\_t2t\_fixtures.py} encode la journée à petit tick comme flux du simulateur~; le test C++ vérifie que le
constructeur d'événements (\omterm{def:ll:the-feed-handler:handler}{feed handler}) la retransforme en les mêmes 6\,000 événements.
\item \textbf{Hors ligne.} Les tests C++ et Rust rejouent les événements à travers le chemin et obtiennent les ordres validés.
\item \textbf{En direct.} \texttt{firm\_ticktotrade\_live\_test.cpp} fait tourner une fois les quatre fils d'exécution et vérifie les ordres, l'ordre des
horodatages et l'absence d'allocation~; \texttt{python bench\_t2t.py} exécute dix passes épinglées et écrit les étapes, les
distributions et le rapport de budget.
\end{tutsteps}
\textbf{Pour aller plus loin.} Remettre les événements d'un paquet au moteur en un seul lot et mesurer à nouveau les étapes de décodage et d'anneau~;
épingler deux fils sur les deux hyperthreads d'un même cœur et observer les étapes qui le partagent.
\end{tutorial}

\section{Construction~: tick-to-trade}\label{bld:ll:build-tick-to-trade-measured:ticktotrade}

\begin{build}
\bfield{Rôle} Les composants du livre assemblés en un chemin de trading unique avec son harnais de mesure~: la preuve qu'ils calculent
ensemble les bons ordres, et la mesure de ce que coûte chacun.
\bfield{Interface} C++20 \texttt{firm::t2t}\,~: \texttt{Path(tick)} avec \texttt{on\_event(event, now, stamp, send)}, \texttt{out},
\texttt{order\_hash}, \texttt{read\_events}, \texttt{Stamp}\,~; \texttt{run\_live(Config)} renvoyant \texttt{Result} avec les
horodatages des ordres, les ordres, les compteurs et les allocations après la phase d'échauffement. Rust \texttt{firm\_ticktotrade}\,~: \texttt{Path},
\texttt{events\_from\_line}, \texttt{order\_hash}. Python\,~: \texttt{make\_t2t\_fixtures.py}~; le \texttt{ll\_t2t.py} du chapitre
calcule les étapes et vérifie le budget avec \texttt{firm.latbudget}.
\bfield{Règles} Le chemin en direct et le rejeu hors ligne produisent les mêmes ordres~; chaque horodatage est pris hors de toute boucle, dans l'ordre du chemin~;
aucune allocation après la phase d'échauffement~; les étapes sont attribuées au premier ordre de chaque événement.
\bfield{Tests d'acceptation} \texttt{code/firm/ticktotrade/}\,~: la ligne enregistrée décodée en événements du \omterm{def:ll:the-order-book-builder:builder}{constructeur de carnet}~; les ordres
hors ligne égaux à l'empreinte validée en C++ et en Rust, sans allocation lors d'une seconde passe~; une exécution en direct avec les mêmes ordres,
les horodatages de chaque ordre dans l'ordre et reçus par la plateforme, et aucune allocation après la phase d'échauffement.
\bfield{Pour aller plus loin} Le serveur de simulation en direct (livre~10) comme plateforme, avec les accusés de réception passant par la machine à états de la passerelle~;
une remise par lots depuis le fil du flux~; des \omterm{def:ll:measuring-latency:hwts}{horodatages matériels} issus d'une carte réseau~; le harnais sur un serveur réglé.
\end{build}

\begin{omsources}
\item Xilinx (AMD), \emph{ef\_vi User Guide}, SF-114063-CD, issue 10, 2019.
\end{omsources}

\section{Exercices}

\begin{exercise}[$\star$]\label{exo:ll:build-tick-to-trade-measured:1}
En utilisant les médianes de la \cref{fig:ll:build-tick-to-trade-measured:stages}, quelle part du temps passé dans le processus est
consommée par les composants propres au livre (livre, stratégie, risque, passerelle)~?
\end{exercise}

\begin{exercise}[$\star$]\label{exo:ll:build-tick-to-trade-measured:2}
Pourquoi l'horloge lue à chaque \omterm{def:ll:build-tick-to-trade-measured:points}{point d'instrumentation} doit-elle être le même compteur sur chaque cœur, et que se passerait-il avec
une horloge par fil d'exécution~?
\end{exercise}

\begin{exercise}[$\star$]\label{exo:ll:build-tick-to-trade-measured:3}
Cinq événements arrivent dans un même paquet et chacun prend environ \qty{100}{\nano\second} pour être décodé et publié. Quelle est, en moyenne, la durée de
l'étape de décodage vécue par les cinq événements, et pourquoi l'étape d'anneau montre-t-elle le même effet~?
\end{exercise}

\begin{exercise}[$\star\star$]\label{exo:ll:build-tick-to-trade-measured:4}
Un événement déclenche deux ordres, chacun envoyé en environ \qty{4}{\micro\second}. Si les étapes étaient attribuées à chaque ordre, quelle étape
absorberait l'attente du second, et de combien sa médiane se déplacerait-elle~?
\end{exercise}

\begin{exercise}[$\star\star$]\label{exo:ll:build-tick-to-trade-measured:5}
Recalculer la projection de la dernière section si l'attente dans le décodage et dans l'anneau était entièrement supprimée. Le chemin projeté
respecte-t-il le budget de trois microsecondes~?
\end{exercise}

\begin{exercise}[$\star\star$]\label{exo:ll:build-tick-to-trade-measured:6}
L'étape de réception inclut le \texttt{sendto} du fil d'échange. Pourquoi ce harnais ne peut-il pas séparer l'appel système de l'émetteur de celui du récepteur, et que changerait
un \omterm{def:ll:measuring-latency:hwts}{horodatage matériel} sur le fil~?
\end{exercise}

\begin{exercise}[$\star\star\star$]\label{exo:ll:build-tick-to-trade-measured:7}
\emph{Programmation.} Transmettre les événements d'un paquet au moteur sous la forme d'un seul message d'anneau et exécuter le moteur sur le lot. Vérifier que les ordres sont inchangés, puis
mesurer à nouveau les étapes de décodage et d'anneau.
\end{exercise}

\begin{exercise}[$\star\star\star$]\label{exo:ll:build-tick-to-trade-measured:8}
\emph{Trouver la faille.} «~Notre tick-to-trade est de \qty{800}{\nano\second}~: nous chronométrons depuis le moment où notre code reçoit le paquet jusqu'au
moment où nous appelons \texttt{send}, et nous publions la moyenne.~»
\end{exercise}

\section{Problème~: Le budget, démontré}

\begin{problem}[{Problème du week-end --- le chemin du livre, étape par étape}]\label{pb:ll:build-tick-to-trade-measured:1}
Utiliser les \cref{fig:ll:build-tick-to-trade-measured:stages,tab:ll:build-tick-to-trade-measured:budget} (\texttt{measured\_stages.csv},
\texttt{measured\_budget.csv}).

\medskip\noindent\textbf{Partie I --- La mesure.}
\begin{enumerate}
\item Énumérer les huit étapes et les \omterm{def:ll:build-tick-to-trade-measured:points}{points d'instrumentation} qui bornent chacune d'elles.
\item Donner les médianes, les 99\up{e} et 99.9\up{e}~centiles de la \omterm{def:ll:where-latency-comes-from:t2t}{latence tick-to-trade} et de la latence dans le processus.
\item Quelle étape a la plus grande médiane, et laquelle le plus grand 99.9\up{e}~centile~?
\item Pourquoi le 99.9\up{e}~centile de l'étape d'anneau est-il si éloigné de sa médiane sur cet ordinateur portable~?
\end{enumerate}

\medskip\noindent\textbf{Partie II --- La vérification.}
\begin{enumerate}[resume]
\item Quelles étapes respectent le budget du chapitre~1 à la médiane, et lesquelles le dépassent, et de quel facteur~?
\item Pourquoi le chiffre de bout en bout du budget est-il composé à partir des étapes des premiers ordres plutôt que tiré de la
distribution du tick-to-trade~?
\item Pourquoi les centiles des étapes ne s'additionnent-ils pas pour donner le centile du total~?
\item Qu'établit la vérification que les tests fonctionnels n'établissent pas~?
\end{enumerate}

\medskip\noindent\textbf{Partie III --- La projection.}
\begin{enumerate}[resume]
\item Remplacer les étapes de réception et d'émission par les chiffres de \omterm{def:ll:linux-tuning:polling}{contournement du noyau} du budget, puis calculer la médiane projetée.
\item Quelle part de l'écart restant relève de l'attente, et quelle part du calcul~?
\item Que changeraient des cœurs isolés sur la figure, et que ne changeraient-ils pas~?
\item Pourquoi la projection n'est-elle pas une mesure, et que faudrait-il pour qu'elle le devienne~?
\end{enumerate}

\medskip\noindent\textbf{Partie IV --- Le verdict.}
\begin{enumerate}[resume]
\item Énoncer le \emph{résultat nommé}~: la médiane mesurée, les 99\up{e} et 99.9\up{e}~centiles par étape et de bout en bout sur cet ordinateur portable,
et la médiane projetée pour un serveur réglé.
\item Quelle part de la médiane dans le processus est le code écrit dans ce livre, et quelle part est de l'attente~?
\item La technique de quel chapitre répond à chacune des trois plus grandes contributions~?
\item Pourquoi les chemins hors ligne et en direct doivent-ils produire les mêmes ordres avant que toute mesure de temps ne soit crédible~?
\item Pourquoi l'absence d'allocation est-elle vérifiée à chaque exécution~?
\item Que faudrait-il mesurer en premier sur un serveur de production~?
\item Quel résultat du livre inspirerait le plus de méfiance si ce harnais n'avait pas été construit~?
\item En une phrase~: que faut-il pour connaître une latence, plutôt que pour y croire~?
\end{enumerate}
\end{problem}

\section{Questions d'entretien}

\begin{interviewq}[$\star$ \iqroles{développeur}]\label{iq:ll:build-tick-to-trade-measured:1}
Qu'est-ce que la \omterm{def:ll:where-latency-comes-from:t2t}{latence tick-to-trade}, et où commencez-vous et arrêtez-vous exactement l'horloge~?
\end{interviewq}

\begin{interviewq}[$\star\star$ \iqroles{développeur}]\label{iq:ll:build-tick-to-trade-measured:2}
Comment instrumenteriez-vous un chemin de trading pour trouver l'étape lente, sans le ralentir~?
\end{interviewq}

\begin{interviewq}[$\star\star$ \iqroles{développeur}]\label{iq:ll:build-tick-to-trade-measured:3}
Votre tick-to-trade médian est correct, mais votre 99.9\up{e}~centile est mille fois plus grand. Où regardez-vous~?
\end{interviewq}

\begin{interviewq}[$\star\star$ \iqroles{développeur}]\label{iq:ll:build-tick-to-trade-measured:4}
Décrivez le parcours d'un paquet depuis la carte réseau jusqu'à votre stratégie et en sens inverse, avec une latence pour chaque étape sur un
serveur réglé.
\end{interviewq}

\begin{interviewq}[$\star\star$ \iqroles{développeur, trader}]\label{iq:ll:build-tick-to-trade-measured:5}
L'équipe stratégie demande un \omterm{def:ll:where-latency-comes-from:budget}{budget de latence}. Comment l'écrivez-vous, et comment prouvez-vous qu'il est respecté~?
\end{interviewq}

\begin{interviewq}[$\star\star\star$ \iqroles{développeur}]\label{iq:ll:build-tick-to-trade-measured:6}
Concevez le harnais de mesure d'un système de trading en production~: que horodate-t-on, où, avec quelle horloge, comment les horodatages
sont-ils rapprochés, et comment les résultats parviennent-ils aux personnes qui en ont besoin~?
\end{interviewq}
